\section{第二部分：白盒缺席时的任务}\label{sec:node}

\subsection{第一部分的结论依赖一个前提}

第一部分把白盒与黑盒放在同一张表上，结论是白盒三轴全胜。但这个结论有一个前提：\textbf{最短路有精确算法}，白盒可以把正确性写进结构里。

节点分类没有这样的算法。没有人能写出“给定图，输出每个节点的类”的正确程序——它是经验性的、依赖数据的。这正是黑盒的主场，也是“可解释性 vs. 性能”真正开始的地方。

于是问题变成：

\begin{quote}
\emph{当白盒缺席时，我们还能不能验证一个解释是对的？}
\end{quote}

这个问题比“解释得好不好”更基本。第一部分已经证明：在最短路任务上，Fidelity$^+$ 高达 0.978 的模型，其预测与真值的相关系数只有 $-0.10$——保真度会被“常数预测器”骗过。如果连保真度都不可信，那么在\emph{没有真值}的任务上，我们凭什么说一个解释是好的？

本文的做法是：\textbf{造一个我们知道答案的环境}。用合成图让标签由社区决定，于是“哪些边携带标签信号”是已知的；再用\emph{消融}把“每条边有多重要”测出来，作为真值。有了真值，才能回答“解释器有没有找到重要的边”。

\subsection{数据：Cora}

Cora 共 2708 个节点、1433 维词袋特征、7 个类别。我们把 Planetoid~\cite{sen2008collective} 的边表对称化并去重：原始记录 10858 条有向边，除以 2 是文献常引的 5429，但其中 151 对节点被重复记录，去重后是 \textbf{5278} 条。同配率（两端都已知标签的边中标签相同的比例）为 \textbf{0.811}。

采用 Planetoid 的公开划分：训练 140（节点 0--139）、验证 500、测试 1000。这不是任意切分——\texttt{ind.cora.x} 恰好就是这 140 个训练节点的特征，可用 \texttt{x == allx[:140]} 断言。

\subsection{模型：GCN 与它的无结构对照}

\textbf{GCN}~\cite{kipf2017gcn}：两层，$\hat A=D^{-1/2}(A+I)D^{-1/2}$，$H^{(l+1)}=\sigma(\hat A H^{(l)}W^{(l)})$，hidden 16，dropout 0.5，Adam（lr $10^{-2}$，weight decay $5\times10^{-4}$），200 轮。

\textbf{MLP 对照}：结构与参数量完全相同（都是 23{,}040），唯一差别是每层不乘 $\hat A$。因此两者的准确率差异只能归因于图结构。

\begin{table}[H]
\centering
\caption{Cora 公开划分上的节点分类（8 个种子）}
\label{tab:cora}
\small
\begin{tabular}{lccc}
\toprule
模型 & 平均准确率 & 最差种子 & 说明 \\
\midrule
GCN & \textbf{0.8162} & 0.8110 & 原论文 81.5\% \\
MLP & 0.5539 & 0.5330 & 无结构基线 \\
\bottomrule
\end{tabular}
\end{table}

GCN 比 MLP 高 26 个点（表~\ref{tab:cora}），说明这个任务上图结构确实提供了 MLP 拿不到的信息。

\subsection{解释器：边掩码与梯度显著性}

\textbf{边掩码解释器}：冻结模型，为目标节点学一张边权重掩码 $m\in[0,1]^{|E|}$：
\begin{equation}
\min_m\ -\log p\bigl(c\mid A\odot m\bigr)+\lambda\,\mathrm{mean}(m),
\end{equation}
其中 $c$ 取模型的全局预测类，$\lambda=0.05$，优化 200 步，温度固定为 1（不做退火——退火会让掩码塌成硬 0/1，产生大量并列）。

实现上有两个必须写清的细节。其一，掩码只在目标节点的 2 跳子图上优化，子图内重新归一化。其二，\textbf{子图的归一化算子必须与训练时完全一致}：度数只累加一次、自环贡献 $d^2$、每条无向边必须双向传入。我们最初写错了这三处中的两处，导致掩码塌成常数、解释完全失效；修复后与训练算子逐元素一致（全图 max 差 $0.0$，删边后 $3\times10^{-8}$）。

\textbf{梯度显著性}：作为对照，也用 $\mathrm{score}_e=\partial\log p(c)/\partial w_e$ 给边排序。它连续、几乎不会并列，实现更简单，但同样依赖正确的归一化算子。

\subsection{攻击：随机加边与 FGA}

\textbf{随机加边}：均匀随机挑当前不存在的边加入，作为基线。

\textbf{FGA}~\cite{zugner2018fga}：把邻接矩阵当连续变量，算 $\partial\mathcal L/\partial A$，按\textbf{带符号}梯度挑“加上之后损失涨得最多”的若干条边加入，每轮重新归一化，迭代 5 轮。

这里有一个容易忽略的方向性问题：若对梯度取绝对值，会同时选中“最伤害”与“最补”的边，净效果反而提升准确率（我们实测把 0.808 抬到 0.875）。必须用带符号梯度。
